% ============================================================
%  HIS Agent — Slide bảo vệ đồ án (theme HUST)
%  Biên dịch: pdflatex main.tex  (chạy 2 lần cho số trang)
%  CẦN: ảnh nền theme trong images/  (hust-red-169-title.png, -section.png,
%       -content-full.png, -end.png). Đổi 169<->43 ở documentclass VÀ \usetheme.
%  Placeholder \imgph{...} = chỗ chèn screenshot: thay bằng
%       \includegraphics[width=\linewidth]{images/ten-anh.png}
% ============================================================
\documentclass[aspectratio=169,t]{beamer}

\usepackage[utf8]{inputenc}
\usepackage[T5]{fontenc}
\usepackage[vietnamese]{babel}
\usepackage{lmodern}
\usepackage{tikz}
\usetikzlibrary{positioning, arrows.meta}
\usepackage[most]{tcolorbox}

\usetheme[169,red]{HUST}

% Thẻ persona (card có viền) — dùng ở slide Người dùng
\tcbset{personacard/.style={
    colframe=hustprimary, colback=hustprimary!4, boxrule=1pt, arc=3pt,
    fonttitle=\bfseries\small, fontupper=\small, coltitle=white,
    left=4pt, right=4pt, top=3pt, bottom=3pt,
    equal height group=persona}}
\showPageNumber

% Khung placeholder cho ảnh chụp (thay bằng \includegraphics khi có ảnh)
\newcommand{\imgph}[2][0.72]{%
  \begin{center}
  \fcolorbox{black!30}{black!5}{%
    \parbox[c][3cm][c]{#1\linewidth}{\centering\color{black!45}\small [Ảnh chụp: #2]}%
  }%
  \end{center}%
}

% ---- Thông tin đề tài ----
\title{HIS Agent}
\subtitle{Trợ lý AI hội thoại cho hệ thống thông tin bệnh viện}
\author{Trương Công Thái Đức}
\institute{Đại học Bách Khoa Hà Nội\\GVHD: PGS.TS. Nguyễn Hồng Quang}

\begin{document}

% ================= BÌA =================
\husttitlepage

% ================= MỤC LỤC (slide 2) =================
\begin{frame}{Nội dung trình bày}
  \vspace{0.4cm}
  \tableofcontents
\end{frame}

% ================= A. ĐẶT VẤN ĐỀ =================
\section{Đặt vấn đề}
\hustsectionpage{Đặt vấn đề}

% --- Người dùng & bối cảnh ---
\begin{frame}{Người dùng \& bối cảnh}
  \begin{center}\small\textit{Bệnh viện đa khoa quy mô vừa — quá tải, đông bệnh nhân}\end{center}
  \vspace{0.1cm}
  \begin{columns}[t]
    \column{0.49\textwidth}
      \begin{tcolorbox}[personacard, title={Bệnh nhân \normalfont\itshape\small— đi khám nhưng bối rối}]
        \begin{itemize}\setlength\itemsep{3pt}
          \item Không biết triệu chứng khám \textbf{khoa nào} → đăng ký nhầm, chờ lâu
          \item Đặt lịch phải gọi/đến trong giờ hành chính, xếp hàng
          \item Triệu chứng đáng lo (sốt, mẩn đỏ) — không rõ \textbf{có nguy hiểm} không
          \item Tra thuốc/dịch trên mạng → nguồn không đáng tin
        \end{itemize}
      \end{tcolorbox}
    \column{0.49\textwidth}
      \begin{tcolorbox}[personacard, title={Bác sĩ \normalfont\itshape\small— quá tải, dễ sai sót}]
        \begin{itemize}\setlength\itemsep{3pt}
          \item Tra hồ sơ rời rạc (tiền sử, xét nghiệm, đơn cũ) → mất thời gian
          \item Nguy cơ \textbf{kê nhầm thuốc} dị ứng (phản ứng chéo khó nhớ)
          \item Nhập tay kết quả xét nghiệm (CBC) → chậm, dễ sai số
          \item Nhớ ngưỡng/phác đồ sàng lọc SXHD → dễ bỏ sót cảnh báo
        \end{itemize}
      \end{tcolorbox}
  \end{columns}
\end{frame}

% --- Bài toán ---
\begin{frame}{Bài toán}
  \begin{itemize}
    \item \textbf{HIS truyền thống}: giao diện form cứng, nhiều bước, \textbf{không hiểu ngôn ngữ tự nhiên}.
    \item \textbf{LLM thuần}: hiểu ngôn ngữ nhưng \textbf{ảo giác}, bịa số/liều → \textbf{không an toàn cho y tế}, không nối dữ liệu bệnh viện.
  \end{itemize}
  \vspace{0.3cm}
  \begin{center}
    \color{hustprimary}\large\textbf{Khoảng trống: cần trợ lý vừa hiểu ngôn ngữ tự nhiên,\\ vừa AN TOÀN và nối được dữ liệu chuẩn.}
  \end{center}
\end{frame}

% --- Mục tiêu ---
\begin{frame}{Mục tiêu}
  Xây dựng \textbf{trợ lý AI hội thoại tiếng Việt} cho bệnh viện:
  \begin{enumerate}
    \item Hiểu \textbf{ngôn ngữ tự nhiên} tiếng Việt cho cả bệnh nhân và bác sĩ.
    \item Quyết định lâm sàng dựa trên \textbf{luật có nguồn} (QĐ 2760, Stockley's) — \textbf{chống ảo giác}.
    \item \textbf{Phân quyền} an toàn theo vai trò (bệnh nhân / bác sĩ).
    \item Lưu trữ theo \textbf{chuẩn FHIR R5} — chuẩn hoá \& liên thông.
  \end{enumerate}
\end{frame}

% ================= B. GIẢI PHÁP & CÔNG NGHỆ =================
\section{Giải pháp \& Công nghệ}
\hustsectionpage{Giải pháp \& Công nghệ}

% --- Kiến trúc 4 lớp ---
\begin{frame}{Tổng quan giải pháp: kiến trúc 4 lớp}
  \centering
  \begin{tikzpicture}[
    box/.style={draw, rounded corners, fill=black!5, minimum width=0.82\textwidth,
                 minimum height=0.85cm, align=center, font=\small},
    arr/.style={-{Latex[length=2mm]}, thick}]
    \node[box] (ui) {Lớp giao diện: React + TypeScript (bệnh nhân / bác sĩ)};
    \node[box, below=0.35cm of ui] (api) {Lớp API: FastAPI (chat, fhir, auth, appointments)};
    \node[box, below=0.35cm of api] (ag) {Lớp Agent: LLM + Tool-Calling Loop};
    \node[box, fill=hustprimary!12, below=0.35cm of ag] (db) {Lớp dữ liệu: Aidbox FHIR R5 \quad|\quad SQLite (tài khoản)};
    \draw[arr] (ui) -- (api); \draw[arr] (api) -- (ag); \draw[arr] (ag) -- (db);
  \end{tikzpicture}
  \vspace{0.2cm}
  \begin{center}\small\color{hustprimary}\textbf{Agent điều phối — Luật quyết định. Mỗi lớp một trách nhiệm.}\end{center}
\end{frame}

% --- Bảng ánh xạ Nhu cầu -> Chức năng ---
\begin{frame}{Từ nhu cầu đến chức năng}
  \footnotesize
  \begin{center}
  \begin{tabular}{@{}p{2.2cm} p{4.6cm} p{4.6cm}@{}}
    \textbf{Người dùng} & \textbf{Nhu cầu} & \textbf{Chức năng đáp ứng} \\[2pt]
    \hline\\[-6pt]
    Bệnh nhân & Không biết khám khoa nào & Gợi ý khoa (RAG) \\
    Bệnh nhân & Đặt lịch phiền & Đặt lịch hội thoại \\
    Bệnh nhân & Lo triệu chứng, sợ SXH & Sàng lọc SXHD (QĐ 2760) \\
    Bệnh nhân & Tra thuốc/dịch không đáng tin & Tra thuốc + số ca dịch (có nguồn) \\
    Bác sĩ & Dễ kê nhầm thuốc dị ứng & Cảnh báo dị ứng 3 tầng \\
    Bác sĩ & Nhập phiếu XN thủ công & Đọc CBC từ ảnh (OCR) \\
    Bác sĩ & Hồ sơ rời rạc & Tóm tắt hồ sơ + kết thúc khám \\
  \end{tabular}
  \end{center}
\end{frame}

% --- Lý do FHIR R5 + Aidbox ---
\begin{frame}{Nền tảng dữ liệu: vì sao FHIR R5 + Aidbox}
  \small
  \begin{columns}[t]
    \column{0.48\textwidth}
      \textbf{\color{hustprimary}FHIR R5 — chọn CHUẨN}
      \begin{itemize}
        \item Chuẩn trao đổi dữ liệu y tế quốc tế (HL7)
        \item Có sẵn mô hình lâm sàng (Patient, Observation…) → không tự thiết kế schema
        \item R5 = bản chính thức mới nhất; prototype mới → không cần hợp R4
      \end{itemize}
    \column{0.48\textwidth}
      \textbf{\color{hustprimary}Aidbox — chọn CÔNG CỤ}
      \begin{itemize}
        \item FHIR server hiện thực chuẩn — không tự code
        \item Hỗ trợ R5 đầy đủ; REST API + search + admin
        \item Dưới là PostgreSQL → vẫn query SQL được
      \end{itemize}
  \end{columns}
  \vspace{0.25cm}
  \begin{center}\color{hustprimary}\textbf{→ Chuẩn hoá + Liên thông; Aidbox lo hạ tầng để tập trung vào AI agent.}\end{center}
\end{frame}

% --- Cách tiếp cận Agent + Tool-calling ---
\begin{frame}{Cách tiếp cận: Agent + Tool-Calling}
  \small
  \begin{itemize}
    \item LLM \textbf{tự chọn công cụ} (tool) để hoàn thành yêu cầu — không theo kịch bản cứng.
    \item Vòng lặp: \emph{LLM đòi tool → hệ thực thi → trả kết quả → LLM tổng hợp} (lặp tới khi đủ).
    \item Mỗi tool là 1 module Python độc lập → thêm chức năng = thêm file (Open/Closed).
    \item \textbf{Quyết định lâm sàng nằm trong tool (luật)}; LLM chỉ điều phối.
  \end{itemize}
  \vspace{0.2cm}
  \begin{center}\color{hustprimary}\textbf{Linh hoạt như hội thoại, an toàn như luật.}\end{center}
\end{frame}

% ================= C. CHỨC NĂNG =================
\section{Chức năng chính}
\hustsectionpage{Chức năng chính}

% --- BN: Gợi ý khoa + Đặt lịch ---
\begin{frame}{Bệnh nhân: Gợi ý khoa (RAG) + Đặt lịch}
  \small
  \begin{columns}[c]
    \column{0.5\textwidth}
      \begin{itemize}
        \item Tả triệu chứng tự nhiên → \textbf{RAG gợi ý đúng khoa} (dù không gõ tên khoa)
        \item Đăng nhập → agent tự đặt lịch: \emph{chọn slot → tạo lịch}, không hỏi lại tên
        \item Lịch ở trạng thái \emph{chờ bác sĩ duyệt}
      \end{itemize}
    \column{0.5\textwidth}
      \imgph[1.0]{Chat gợi ý khoa + đặt lịch}
  \end{columns}
\end{frame}

% --- BN: Sàng lọc SXHD ---
\begin{frame}{Bệnh nhân: Sàng lọc Sốt xuất huyết}
  \small
  \begin{columns}[c]
    \column{0.5\textwidth}
      \begin{itemize}
        \item Nhận diện triệu chứng → hiện \textbf{phiếu tích} dịch tễ / triệu chứng / dấu hiệu cảnh báo
        \item Phân độ bằng \textbf{luật theo QĐ 2760/QĐ-BYT} — không phải LLM tự phán
        \item Khuyến cáo + trích nguồn; cảnh báo → đề xuất đặt lịch khẩn
      \end{itemize}
    \column{0.5\textwidth}
      \imgph[1.0]{Phiếu SXHD + kết quả phân độ}
  \end{columns}
\end{frame}

% --- BN: Tra thuốc + số ca dịch ---
\begin{frame}{Bệnh nhân: Tra thuốc \& Số ca dịch bệnh}
  \small
  \begin{columns}[c]
    \column{0.5\textwidth}
      \begin{itemize}
        \item \textbf{Tra thuốc}: RAG tra NCBI (PubMed/StatPearls), mỗi ý kèm \textbf{trích dẫn [n]} → chống bịa liều
        \item \textbf{Số ca dịch}: lấy số thật từ \textbf{HCDC}, có tuần + nguồn; thiếu dữ liệu thì báo, không bịa
      \end{itemize}
    \column{0.5\textwidth}
      \imgph[1.0]{Trả lời thuốc có [n] nguồn}
  \end{columns}
\end{frame}

% --- BS: Cảnh báo dị ứng 3 tầng ---
\begin{frame}{Bác sĩ: Cảnh báo dị ứng khi kê đơn (3 tầng)}
  \small
  \begin{columns}[c]
    \column{0.46\textwidth}
      Bệnh nhân dị ứng \textbf{Penicillin}:
      \begin{itemize}
        \item \textbf{Tier 1} — banner cảnh báo sớm đầu form
        \item \textbf{Tier 2} — chọn Cefuroxime → cảnh báo \emph{phản ứng chéo}, cần lý do override
        \item \textbf{Tier 3} — chọn Amoxicillin → \emph{chặn kê}
      \end{itemize}
      \textit{Bắt cả phản ứng chéo, không chỉ trùng tên (Stockley's, BNF).}
    \column{0.54\textwidth}
      \imgph[1.0]{Modal kê đơn: banner + dòng chặn/cảnh báo}
  \end{columns}
\end{frame}

% --- BS: CBC ảnh + Kết thúc khám ---
\begin{frame}{Bác sĩ: CBC từ ảnh (OCR) \& Kết thúc khám}
  \small
  \begin{columns}[c]
    \column{0.5\textwidth}
      \begin{itemize}
        \item Upload ảnh phiếu → \textbf{OCR} trích PLT/WBC/HCT → lưu FHIR → tự sàng lọc SXHD
        \item Kết thúc khám: tạo \emph{ClinicalImpression} (chẩn đoán + kế hoạch) + gửi email tóm tắt cho bệnh nhân
      \end{itemize}
    \column{0.5\textwidth}
      \imgph[1.0]{CBC từ ảnh / Kết thúc khám}
  \end{columns}
\end{frame}

% ================= D. KỸ THUẬT NỔI BẬT =================
\section{Điểm kỹ thuật nổi bật}
\hustsectionpage{Điểm kỹ thuật nổi bật}

\begin{frame}{Điểm kỹ thuật cốt lõi}
  \begin{itemize}
    \item \textbf{Agent tool-calling}: LLM tự chuỗi công cụ; thêm chức năng bằng plugin (Open/Closed).
    \item \textbf{2 kiểu RAG}: vector (gợi ý khoa) + tra API có trích dẫn (thông tin thuốc).
    \item \textbf{Luật lâm sàng có nguồn}: QĐ 2760 (SXHD), Stockley's/BNF (dị ứng) → chống ảo giác.
    \item \textbf{OCR-first}: đọc số bằng model chuyên dụng trước, chính xác hơn cho xét nghiệm.
    \item \textbf{Phân quyền 2 lớp}: prompt + code (ép patient\_id từ session) → chống prompt injection.
  \end{itemize}
  \vspace{0.2cm}
  \begin{center}\color{hustprimary}\large\textbf{LLM điều phối — LUẬT quyết định → đáng tin, không ảo giác.}\end{center}
\end{frame}

% ================= E. ĐÁNH GIÁ & KẾT =================
\section{Kiểm thử \& Kết luận}
\hustsectionpage{Kiểm thử \& Kết luận}

% --- Kiểm thử & hiệu năng ---
\begin{frame}{Kiểm thử \& hiệu năng}
  \small
  \begin{itemize}
    \item \textbf{9 luồng chức năng (A–I)}: đạt toàn bộ.
    \item \textbf{Đánh giá nguy cơ SXHD từ CBC}: kết quả khớp ngưỡng QĐ 2760 (Bảng 4.1).
    \item \textbf{Cảnh báo dị ứng 3 tầng}: 5/5 ca đúng kỳ vọng (chặn/cảnh báo/an toàn).
    \item \textbf{Phân quyền}: bệnh nhân không truy được hồ sơ người khác (code ép patient\_id).
    \item \textbf{Thời gian phản hồi}: luật cục bộ <1.5s; gọi API ngoài ~3–4.8s.
  \end{itemize}
\end{frame}

% --- Hạn chế & hướng phát triển ---
\begin{frame}{Hạn chế \& hướng phát triển}
  \small
  \begin{columns}[t]
    \column{0.48\textwidth}
      \textbf{\color{hustprimary}Hạn chế}
      \begin{itemize}
        \item LLM chạy qua API cloud → dữ liệu ra ngoài viện
        \item Mới kiểm thử kỹ thuật, chưa kiểm định lâm sàng
        \item Sàng lọc mới cho SXHD
      \end{itemize}
    \column{0.48\textwidth}
      \textbf{\color{hustprimary}Hướng phát triển}
      \begin{itemize}
        \item LLM tiếng Việt \textbf{on-premise} (bảo mật dữ liệu)
        \item \textbf{Kiểm định lâm sàng} trước triển khai
        \item Mở rộng bệnh (COVID-19, sốt rét, viêm phổi)
        \item Fine-tune y khoa VN, multi-agent, app di động
      \end{itemize}
  \end{columns}
\end{frame}

% --- Kết luận ---
\begin{frame}{Kết luận}
  \begin{itemize}
    \item Xây dựng thành công \textbf{trợ lý AI hội thoại} cho HIS: hỗ trợ bệnh nhân \& bác sĩ.
    \item Điểm cốt lõi: \textbf{hiểu ngôn ngữ tự nhiên} nhưng \textbf{an toàn nhờ luật có nguồn}, chuẩn FHIR R5.
    \item Kiến trúc plugin + phân quyền \textbf{sẵn sàng mở rộng}.
  \end{itemize}
  \vspace{0.3cm}
  \begin{center}\color{hustprimary}\large\textbf{AI hiểu ngôn ngữ — Luật đảm bảo an toàn.}\end{center}
\end{frame}

% ================= CẢM ƠN =================
\hustthankyou

\end{document}
